%% nilHMM methods paper — preprint.
%% BASE: RILAB bioRxiv template (rilabRxiv.cls, genetics.bst) — the lab's
%% preprint styling/workflow.
%% STRUCTURE: Oxford Bioinformatics — structured abstract (Motivation/Results/
%% Availability/Contact), Intro / Materials & Methods / Results / Discussion,
%% and OUP back-matter (Conflicts, Funding, Data availability, Author
%% contributions). Reproduced in RILAB-compatible LaTeX; port to
%% oup-authoring-template.cls at submission (bundle kept in zealhmm/agent/).

\documentclass[9pt,twocolumn,twoside]{rilabRxiv}
% Use the documentclass option 'lineno' to view line numbers
\setlength{\marginparwidth}{2cm}
\usepackage[textsize=tiny,colorinlistoftodos]{todonotes} % comments in margins
\definecolor{cornflowerblue}{rgb}{0.39, 0.58, 0.93}

%%%%%%% Add comments in color
\newcommand{\ms}[1]{{\small \textcolor{green}{#1}}}
\newcommand{\jri}[1]{{\small \textcolor{red}{#1}}}
\newcommand{\citex}[1]{{\small \textcolor{red}{CITE(#1)}}}
\newcommand{\X}{{\textcolor{red}{X}}}

\newcolumntype{b}{X}
\newcolumntype{s}{>{\hsize=.5\hsize}X}

% Set supplement numbers to S and start counting newly
\newcommand{\beginsupplement}{%
        \setcounter{table}{0}
        \renewcommand{\thetable}{S\arabic{table}}%
        \setcounter{figure}{0}
        \renewcommand{\thefigure}{S\arabic{figure}}%
     }

\usepackage{hyperref}

\title{Simulation-calibrated HMM ancestry calling for near-isogenic lines across sequencing modalities (nilHMM)}

\author[$\ast$,1]{Fausto Rodr\'iguez-Zapata}
\author[$\ast$]{Author list to be finalized}
\author[$\ast$,1]{Rub\'en Rell\'an-\'Alvarez}

\affil[$\ast$]{Department of Molecular and Structural Biochemistry, North Carolina State University, Raleigh, NC, USA}

\keywords{ancestry inference, hidden Markov model, near-isogenic lines, simulation calibration, low-coverage sequencing, RNA-seq}

\runningtitle{Simulation-calibrated HMM ancestry calling} % For use in the footer
\runningauthor{Rodr\'iguez-Zapata \textit{et al.}}

%%% Abstract — Bioinformatics structured form %%%%%%%%%%%%%%%%%%
\begin{abstract}
\noindent\textbf{Motivation:} Ancestry calls for near-isogenic lines (NILs) from
low-cost sequencing normally require calibrating a caller against high-coverage
truth samples, of which few are usually available. We ask whether the calibration
can instead come from cheap simulations.\\
\textbf{Results:} We present nilHMM, a unified duration-aware three-state
(REF/HET/ALT) hidden Markov model whose emission is the only swappable axis,
expressing four ancestry callers (nNIL, a Julia-free RTIGER port, a binned
Gaussian caller, and ATLAS) through one engine. We show that the caller
parameters --- rigidity, error rate, and emission means --- can be calibrated from
Broman \texttt{simcross} simulations rather than from scarce high-coverage truth,
and that the calibrated parameters transfer across skim WGS and 3\textquotesingle{}
RNA-seq (BRB-seq) modalities on real BZea teosinte-introgression NILs.\\
\textbf{Availability:} The \texttt{nilHMM} R package
(\url{https://sawers-rellan-labs.github.io/nilhmm/}) and the \texttt{zealhmm}
analysis repository (\url{https://sawers-rellan-labs.github.io/zealhmm/}) are
openly available; simulated benchmark data are regenerated with pinned seeds.\\
\textbf{Contact:} \href{mailto:thegemmalab@ncsu.edu}{thegemmalab@ncsu.edu}\\
\textbf{Supplementary information:} Supplementary data are available at
\textit{Bioinformatics} online.
\end{abstract}
%%%%%%%%%%%%%%%%%%%%%%%%%%

\setboolean{displaycopyright}{true}

\begin{document}

\maketitle
\thispagestyle{firststyle}
\correspondingauthoraffiliation{
Department of Molecular and Structural Biochemistry, North Carolina State
University, Raleigh, NC, USA. E-mail: corresponding@ncsu.edu}
\vspace{-11pt}%

\setboolean{displaylineno}{true}
\ifthenelse{\boolean{displaylineno}}{\linenumbers}{}

%%%%%%%%%%%%%%%%%%%%%%%%%%%%%%%%%%%%%%%%%%%%%%%%%%%%%%
\section{Introduction}
%%%%%%%%%%%%%%%%%%%%%%%%%%%%%%%%%%%%%%%%%%%%%%%%%%%%%%
\lettrine[lines=2]{\color{color2}N}{}ear-isogenic lines (NILs) carry small donor
introgressions on an otherwise recurrent-parent background, and mapping their
ancestry from sequencing underpins introgression and QTL studies. In the BZea
project, teosinte donor segments are introgressed into a B73 background and read
out by several sequencing modalities of differing cost and depth.

\todo[inline]{Motivate ancestry calling from cheap sequencing; the calibration
bottleneck --- Holland's high-coverage-truth calibration vs. the few truth
samples actually available (our target-seq analogue). State scope: DNA (skim WGS)
and RNA (BRB-seq) modalities. Our contribution: a simulation-calibrated unified
caller family.}

%%%%%%%%%%%%%%%%%%%%%%%%%%%%%%%%%%%%%%%%%%%%%%%%%%%%%%
\section{Materials and Methods}
%%%%%%%%%%%%%%%%%%%%%%%%%%%%%%%%%%%%%%%%%%%%%%%%%%%%%%
\label{sec:materials:methods}

\subsection{A unified duration-aware three-state HMM engine}
One HMM over states REF/HET/ALT. Two axes are swappable and define each caller:
the \emph{emission} (count/BetaBinomial vs.\ categorical genotype-confusion) and
the \emph{duration/rigidity} layer (geometric, rigidity, or HSMM). Design
start-priors come from the pedigree (BC2S2/BC2S3). The engine is data-agnostic:
$(\text{data}, \text{params}) \rightarrow \text{calls}$.
\todo{Fig F1: engine + caller-family schematic.}

\subsection{The callers}
\subsubsection{nNIL} Count/BetaBinomial emission with a geometric duration; the
Holland near-isogenic-line caller \citex{Holland nNIL}.
\subsubsection{RTIGER} A Julia-free port in which a rigidity layer replaces the
free geometric duration \citex{RTIGER}.
\subsubsection{binhmm} A three-state Gaussian emission on per-bin allele
frequency with REF anchored at the coverage floor (dense skim only).
\subsubsection{ATLAS} Competitive transcript-level alignment scores $\rightarrow$
hard genotype $\rightarrow$ categorical emission; a GOOGA-shadow \citex{GOOGA},
built for expression-driven RNA depth and allele-specific expression.

\subsection{Data sources and marker sets}
Skim (low-coverage WGS with teosinte--B73 wideseq filtering), BRB-seq
(3\textquotesingle{} RNA-seq, wideseq-thinned), and MolBreeding target-seq (45K,
high coverage; completeness only).

\subsection{The \texttt{simcross} simulation model}
Forward-simulated BC2S3 (and bulked BC2S2 for skim) NILs on the consensus
genetic map with interference, degraded to each source's depth/error regime
(coverage-driven missingness $\text{missing}(\lambda)=\pi+(1-\pi)e^{-k\lambda}$
and per-read error).

\subsection{Calibration and validation protocol}
Fit rigidity $r$ by minimum KS distance to the simulated fragment-size
distribution (not MLE), and $\text{err}$/emission means on degraded-simulation
truth. Validate on held-out simulations: per-bin state accuracy, segment-boundary
error, breakpoint count, donor-footprint ROC, and HET-vs-ALT calibration.

%%%%%%%%%%%%%%%%%%%%%%%%%%%%%%%%%%%%%%%%%%%%%%%%%%%%%%
\section{Results}
%%%%%%%%%%%%%%%%%%%%%%%%%%%%%%%%%%%%%%%%%%%%%%%%%%%%%%
\subsection{Simulation calibration}
Fitted parameters per source $\times$ caller and calibration curves; why $r$ is a
minimum-distance rather than a maximum-likelihood estimate. \todo{Fig F2.}

\subsection{Simulation benchmark and caller ranking}
Accuracy, boundary error, breakpoint count, donor-footprint ROC, and HET/ALT
calibration on simulation truth. \todo{Fig F3.}

\subsection{Application to the paired BRB + skim cohort}
Transfer the simulation-calibrated parameters to the $\sim$400 NILs sequenced by
both modalities; the paired design isolates method from sample. Cross-source and
cross-caller concordance: pairwise segment agreement, donor-footprint Jaccard,
breakpoint concordance, control cleanliness on B73. \todo{Fig F4.} Chromosome
painting appears only as a supplementary QC illustration.

\subsection{B1 anthocyanin mapping (deferred)}
\todo[inline]{Included if the B1 phenotype table and NIL panel are ready in time;
otherwise a future-application paragraph in the Discussion.}

%%%%%%%%%%%%%%%%%%%%%%%%%%%%%%%%%%%%%%%%%%%%%%%%%%%%%%
\section{Discussion}
%%%%%%%%%%%%%%%%%%%%%%%%%%%%%%%%%%%%%%%%%%%%%%%%%%%%%%
\todo[inline]{Simulations can replace truth samples for calibration; limitations
(zygosity identifiability ceiling, ASE for ATLAS, target-seq too sparse to
calibrate from); which caller/source to use when.}

%%%%% Bioinformatics back-matter %%%%%
\section*{Conflicts of interest}
None declared.

\section*{Funding}
\todo{Funding sources.}

\section*{Data availability}
The \texttt{nilHMM} package, the \texttt{zealhmm} analysis repository, and the
simulated benchmark data (regenerated with pinned seeds) are openly available;
see Availability.

\section*{Author contributions}
\todo{CRediT statement once the author list is finalized.}

\section{Acknowledgments}
\todo{Acknowledgments.}

\bibliography{references}

%\pagebreak
\onecolumn
\section*{Supplement}
\beginsupplement
\todo[inline]{Supplementary chromosome-painting sanity check (source $\times$
caller tracks) and any supporting tables.}

\end{document}
